% =====================================================================
%  Выводы по лабораторной работе.
% =====================================================================

\labpart{Выводы}

В ходе выполнения лабораторной работы реализована автоматизированная
система синтеза речи по варианту~1 (английский язык, научные статьи по
computer science) как новый сценарий использования поверх архитектуры,
разработанной в ЛР1--ЛР4. Синтез вынесен в отдельный микросервис
\texttt{speech-service} -- по тому же принципу изоляции тяжёлых
ML-зависимостей, что и определение языка и нейросетевой перевод; базу
данных синтез не затрагивает, существующие контракты системы не
нарушены.

Обязательные возможности реализованы следующим образом. Ввод текста:
поле произвольного текста, кнопка «Озвучить» рядом с документом,
рефератом и переводом, а также кнопка «Озвучить из буфера», которая
читает текст, скопированный в любом другом приложении (озвучивание по
указателю мыши в чужом окне невозможно в браузерной среде по
соображениям безопасности; заменой служит буфер обмена, разрешённый
заданием). Воспроизведение: нейросетевой синтез локальной моделью
Piper, потоковая передача звука и воспроизведение через Web Audio API.
Настройки: 10 английских голосов (7 американских и 3 британских),
темп от $0{,}5\times$ до $2{,}0\times$ и громкость от 0 до 100\,\%,
которые сохраняются между сеансами.

Ключевое решение -- разбиение текста на фрагменты по предложениям с
упреждающей выборкой следующего фрагмента: без него озвучивание текста
в 2298 символов началось бы через 63{,}6~с, а с ним -- через 2{,}9~с,
причём на длинном тексте пауз между фрагментами не возникает вовсе
(синтез идёт в среднем вдвое быстрее звучания, RTF около 0{,}48).
Круговая проверка разборчивости (синтез с последующим распознаванием)
дала 2{,}6\,\% ошибочных слов на предложениях из предметной области,
с ошибками только на редких терминах.

Измерения выявили и слабые места, которые описаны в разделе «Анализ
результатов и предложения по улучшению»: задержку до первого звука в
2{,}7--3{,}5~с, дополнительные 3{,}8~с на первое обращение к каждому
голосу и отсутствие вытеснения загруженных голосов -- после обращения ко
всем десяти голосам и длинных синтезов память контейнера приблизилась
к лимиту (3{,}84 из 4~ГиБ). Для каждого из них предложено
решение, наиболее важное из которых -- кэш голосов с вытеснением.
